\section{Limitações da Pesquisa}
\label{sec:limitations}

As limitações deste estudo são organizadas segundo a taxonomia de validade de \citet{Wohlin2012} e devem ser consideradas na interpretação dos resultados.

\textbf{Validade interna.} A análise é descritiva e associativa, não causal. Não é possível afirmar que a \textit{formação docente} ou a \textit{infraestrutura escolar} causam IDEB mais alto. Municípios com gestão mais eficaz tendem a contratar mais docentes qualificados e a obter melhores resultados ao mesmo tempo, o que torna difícil separar o efeito de cada variável do efeito da qualidade de gestão municipal. A segmentação reduz o confundimento pela renda, mas não o elimina: no \textit{cluster} intermediário, a renda \textit{per capita} ainda difere marginalmente entre os grupos de alto e de baixo IDEB.

\textbf{Validade externa.} Os resultados se referem aos municípios maranhenses, à rede municipal, aos anos iniciais do ensino fundamental e à edição 2025 do IDEB. A transposição para outros estados exige replicar o procedimento, e não apenas os achados. Os 15 municípios de alta renda, que reúnem a capital e os principais polos regionais, não têm amostra suficiente para a comparação por quartis, e os fatores associados ao desempenho nesse grupo permanecem sem investigação.

\textbf{Validade de construto.} As duas variáveis de segmentação descrevem a distribuição de renda e são correlacionadas entre si ($-0{,}87$). Dimensões socioeconômicas não monetárias, como a escolaridade dos adultos e as condições de saúde, não entram na segmentação; a ablação mostra que a pobreza extrema acrescenta informação à renda média, mas a partição ainda é essencialmente ordenada ao longo de um eixo de renda. A titulação é uma medida indireta da qualidade docente, e a presença de quadra, uma medida indireta do investimento em infraestrutura. As variáveis são agregados municipais, e as associações valem para as redes municipais, não para escolas, professores ou alunos \citep{robinson1950}. A formação docente da Sinopse inclui os docentes de todas as redes do município; a verificação feita apenas com a rede municipal, que reúne 93\% a 95\% desses docentes nos grupos analisados, reproduz os resultados. Por fim, o próprio IDEB incorpora a aprovação, o que cria relação parcialmente mecânica com as variáveis de fluxo \citep{soares2013}.

\textbf{Validade de conclusão.} As comparações envolvem grupos de 23 a 34 municípios, o que limita o poder estatístico. Com a correção para comparações múltiplas, apenas a quadra esportiva e a correção de fluxo no \textit{cluster} intermediário mantêm valor $q$ próximo de 0,05. A formação docente é sustentada pela consistência entre os dois \textit{clusters} e entre os dois procedimentos de comparação, e não por um teste isolado, e os resultados devem ser lidos como exploratórios. Os modelos preditivos treinados com as variáveis do estudo não superaram o acaso na separação dos grupos nem a média na previsão do IDEB. O Censo Escolar não captura processos pedagógicos, gestão, clima escolar ou engajamento familiar, que são mediadores importantes do desempenho e provavelmente respondem pela maior parte da variação não explicada.
